\SourcePage{394}
\SGAsection{5}{附录：整性定理}
\begin{center}P. Deligne 著\end{center}

\textbf{5.0。}设 $p$ 为素数，$q$ 为 $p$ 的幂，$\mathbb F_q$ 为有 $q$ 个元素的域，$E$ 为一个\Term{number-field}{数域}在整除某个异于 $p$ 的素数 $\ell$ 的地方处的\Term{completion}{完备化}，$T$ 为一组素数（例如 $T=\varnothing$）。

$E$ 的一个\Term{algebraic-closure}{代数闭包}中的元素 $\alpha$ 称为\Term{T-integral}{$T$-整的}，若它在 $\mathbb Q$ 上\Term{algebra}{代数}，且在 $\mathbb Z[(1/r)_{r\in T}]$ 上整。$E$ 上\Term{finite-dimensional-vector-space}{有限维向量空间}的\Term{endomorphism}{自同态} $F$ 称为\Term{T-integral}{$T$-整的}，若其\Term{eigenvalue}{特征值}均为\Term{T-integral}{$T$-整的}。

设 $k$ 为有 $q'$ 个元素的\Term{finite-field}{有限域}，$\bar k$ 为 $k$ 的一个\Term{algebraic-closure}{代数闭包}，$\mathfrak G=\operatorname{Gal}(\bar k/k)$ 为 $k$ 的\Term{galois-group}{Galois 群}。用 $f\in\mathfrak G$ 表示由
\[
f(x)=x^{q'}
\]
给出的 Frobenius 置换，用 $F$ 表示“几何”Frobenius，即 $f$ 在 $\mathfrak G$ 中的逆。$\mathfrak G$ 在 $E$ 上\Term{finite-dimensional-vector-space}{有限维向量空间}中的一个\Term{representation}{表示} $\rho$ 称为\Term{T-integral}{$T$-整的}，若 $\rho(F)$ 为\Term{T-integral}{$T$-整的}。若 $X$ 是 $\mathbb F_q$ 上\Term{finite-type}{有限型}\Term{scheme}{概形}，$x$ 为 $X$ 的一个\Term{closed-point}{闭点}，则用 $F_x$ 表示 $k(x)$ 的\Term{galois-group}{Galois 群}中的几何 Frobenius 元素；若 $\bar x$ 是 $X$ 的一个位于 $x$ 之上的\Term{geometric-point}{几何点}，$G$ 为 $X$ 上的一个 $E$-\Term{sheaf}{层}，则 $F_x$ 通过结构搬移\Term{action}{作用}于\Term{fibre}{纤维} $G_{\bar x}$。

\begin{SGAresult}{定义}{5.1}
$\mathbb F_q$ 上\Term{finite-type}{有限型}\Term{scheme}{概形} $X$ 上的\Term{constructible}{可构造} $E$-\Term{sheaf}{层} $G$ 称为\Term{T-integral}{$T$-整的}，若对 $X$ 的每个位于\Term{closed-point}{闭点} $x$ 之上的\Term{geometric-point}{几何点} $\bar x$，$F_x$ 都是 $G_{\bar x}$ 的一个\Term{T-integral}{$T$-整}\Term{automorphism}{自同构}。
\end{SGAresult}

\SourcePage{395}
若 $a:X\to\operatorname{Spec}(\mathbb F_q)$ 是 $\mathbb F_q$ 上\Term{finite-type}{有限型}\Term{scheme}{概形}，$G$ 是 $X$ 上的\Term{constructible}{可构造} $E$-\Term{sheaf}{层}，则用 $\underline H^*(X,G)$ 表示 $\operatorname{Spec}(\mathbb F_q)_{\text{ét}}$ 上的 $E$-\Term{sheaf}{层} $R^*a_*(G)$；若 $\overline{\mathbb F}_q$ 为 $\mathbb F_q$ 的一个\Term{algebraic-closure}{代数闭包}，$\bar X$ 为从 $X$ 经 $\mathbb F_q\to\overline{\mathbb F}_q$ 的\Term{scalar-extension}{标量扩张}所得的\Term{scheme}{概形}，则 $\underline H^*(X,G)$ 与\Term{cohomology}{上同调}群 $H^*(\bar X,G)$ 等同，后者看作\Term{galois-module}{Galois 模}，即带有几何 Frobenius $F$ 的\Term{action}{作用}。\Term{proper-support-cohomology}{固有支撑上同调}也采用同样的记法。

\begin{SGAresult}{\Term{lemma}{引理}}{5.2.1}
设 $G$ 为 $\mathbb F_q$ 上\Term{dimension}{维数} $\leq n$ 的\Term{finite-type}{有限型}\Term{scheme}{概形} $X$ 上的\Term{constructible}{可构造} $E$-\Term{sheaf}{层}。

\textup{(i)} 存在 $X$ 的一个零维\Term{closed-subset}{闭子集} $Y$，使显然的映射 $\underline H^0(X,G)\to\underline H^0(Y,G)$ 为\Term{injective}{单射}。

\textup{(ii)} 若 $X$ \Term{separated}{分离}，$U$ 为 $X$ 的一个\Term{open-subscheme}{开子概形}，其补集 $Z$ 的\Term{dimension}{维数} $<n$，则存在 $U$ 的一个零维\Term{closed-subset}{闭子集} $Y$，以及一个\Term{surjective}{满射}
\[
\underline H^0(Y,G)(-n)\longrightarrow\underline H_c^{2n}(X,G).
\]
\end{SGAresult}
\begin{proof}
断言 (i) 是显然的。证明 (ii)。将 $X$ 换成 $X_{\mathrm{red}}$，并缩小 $U$ 后，可以假定 $U$ \Term{smooth}{光滑}，且 $G|U$ \Term{twisted-constant}{扭常值}。因为 $\dim Z<n$，有
\[
0=\underline H_c^{2n-1}(Z,G)\longrightarrow
\underline H_c^{2n}(U,G)\xrightarrow{\sim}
\underline H_c^{2n}(X,G)\longrightarrow
\underline H_c^{2n}(Z,G)=0,
\]
由\Term{poincare-duality}{Poincaré 对偶}（SGA 4 第 XVIII 讲及 [2]），又有
\[
\underline H_c^{2n}(U,G)\simeq\underline H^0(U,G^\vee(n))^\vee.
\]
于是由 (i) 得到结论。
\end{proof}

\SourcePage{396}
\begin{SGAresult}{\Term{theorem}{定理}}{5.2.2}
设 $X$ 为 $\mathbb F_q$ 上\Term{dimension}{维数} $\leq n$ 的\Term{separated}{分离}\Term{finite-type}{有限型}\Term{scheme}{概形}，$G$ 为\Term{T-integral}{$T$-整}\Term{constructible}{可构造} $E$-\Term{sheaf}{层}。则 $\underline H_c^i(X,G)$ 的\Term{endomorphism}{自同态} $F$ 和 $q^{n-i}F$ 均为\Term{T-integral}{$T$-整的}。
\end{SGAresult}
\begin{proof}
$\dim X\leq0$ 的情形是平凡的，而\Term{theorem}{定理}中 $i=0$ 和 $i=2n$ 的情形由 5.2.1 得到。设 $Z(X,G,t)$ 为\Term{formal-power-series}{形式幂级数}
\[
Z(X,G,t)=\prod_{x\in X}\det\bigl(1-F_xt^{\deg(x)};G_{\bar x}\bigr)^{-1}\in E[[t]].
\]
由 SGA 5 第 XV 讲 3.2，有
\[
Z(X,G,t)=\prod_i\det\bigl(1-Ft;\underline H_c^i(X,G)\bigr)^{(-1)^{i+1}}.
\tag{5.2.3}
\]
先对 $n=1$ 证明 5.2.2。此时，当 $i>2$ 时，$\underline H_c^i(X,G)=0$（SGA 4 第 X 讲 4.3 或第 XVII 讲 5.2.8.1）。由假设，\Term{formal-power-series}{形式幂级数} $Z(X,G,t)$ 的系数均为\Term{T-integral}{$T$-整的}。根据 5.2.2 中 $i=0$ 和 $i=2n=2$ 的情形，$Z(X,G,t)$ 与
\[
\det\bigl(1-Ft;\underline H_c^0(X,G)\bigr)\cdot
\det\bigl(1-Ft;\underline H_c^2(X,G)\bigr)
\]
的乘积仍具有\Term{T-integral}{$T$-整}系数。这个乘积正是 $\det(1-Ft;\underline H_c^1(X,G))$，由此得到 $n=1$ 时的\Term{theorem}{定理}。

现对 $n=\dim X$ 归纳。设 $U$ 为 $X$ 的一个 Zariski \Term{open-subscheme}{开子概形}，其补集 $Z$ 的\Term{dimension}{维数} $<n$，并且存在一个\Term{morphism}{态射} $f:U\to Y$，其目标为一条\Term{curve}{曲线}，\Term{fibre}{纤维}的\Term{dimension}{维数} $\leq n-1$。
\SourcePage{397}
\[
\begin{tikzcd}
U\arrow[r,hook,"j"]\arrow[d,"f"']&X&Z\arrow[l,hook',"i"']\\
Y&&
\end{tikzcd}
\]
\Term{exact-sequence}{正合列}
\[
0\longrightarrow j_!j^*G\longrightarrow G\longrightarrow i_!i^*G\longrightarrow0
\]
诱导\Term{exact-sequence}{正合列}
\[
\underline H_c^i(U,G)\longrightarrow\underline H_c^i(X,G)\longrightarrow\underline H_c^i(Z,G).
\]
所以由归纳，只须对 $U$ 上的\Term{sheaf}{层} $G$ 验证\Term{theorem}{定理}。将归纳假设用于 $f$ 的\Term{fibre}{纤维}，并利用\Term{proper-base-change}{固有基变换定理}，可知 $Y$ 上的\Term{sheaf}{层} $R^if_!G$ 为\Term{T-integral}{$T$-整的}，\Term{sheaf}{层} $R^if_!G(i-(n-1))$ 也为\Term{T-integral}{$T$-整的}。$f$ 的 Leray \Term{spectral-sequence}{谱序列}为
\[
E_2^{pq}=\underline H_c^p(Y,R^qf_!G)\Longrightarrow\underline H_c^{p+q}(U,G).
\]
由上述结果及已经对 $Y$ 证明的\Term{theorem}{定理}，$E_2^{pq}$ 为\Term{T-integral}{$T$-整的}，$E_2^{pq}(p+q-n)$ 也为\Term{T-integral}{$T$-整的}。\Term{theorem}{定理}随之得证。
\end{proof}

\begin{SGAresult}{推论}{5.3}
设 $X$ 为 $\mathbb F_q$ 上\Term{dimension}{维数} $\leq1$ 的\Term{finite-type}{有限型}\Term{scheme}{概形}，$U$ 为 $X$ 的一个\Term{open-subscheme}{开子概形}，$j:U\hookrightarrow X$ 为包含\Term{morphism}{态射}，$G$ 为 $U$ 上的\Term{constructible}{可构造} $E$-\Term{sheaf}{层}。则：
\SourcePage{398}

\textup{(a)} 若 $G$ 为\Term{T-integral}{$T$-整的}，则 $j_*G$ 为\Term{T-integral}{$T$-整的}。

\textup{(b)} 若对每个 $x\in U$，$G_{\bar x}^\vee$ 均为\Term{T-integral}{$T$-整的}，则对每个 $x\in X$，$(j_*G)_{\bar x}^\vee$ 均为\Term{T-integral}{$T$-整的}。
\end{SGAresult}
\begin{proof}
先归约到 $X$ \Term{affine}{仿射}的情形，再归约到 $X$ \Term{projective}{射影}的情形。

用\Term{exact-sequence}{正合列}
\[
0\longrightarrow j_!G\longrightarrow j_*G\longrightarrow H\longrightarrow0
\]
定义 $H$。若 $Y$ 为 $H$ 的（\Term{finite}{有限}）\Term{support}{支撑}，则相应的\Term{cohomology}{上同调}\Term{exact-sequence}{正合列}给出\Term{exact-sequence}{正合列}
\[
\underline H^0(U,G)\longrightarrow\underline H^0(Y,H)\longrightarrow\underline H_c^1(U,G).
\]
由 5.2.1 和\Term{theorem}{定理} 5.2，$\underline H^0(U,G)$ 与 $\underline H_c^1(U,G)$ 为\Term{T-integral}{$T$-整的}；故 $\underline H^0(Y,H)$ 也为\Term{T-integral}{$T$-整的}，这就证明了 (a)。

若 $X'$ 为 $X_{\mathrm{red}}$ 的纯一维部分的\Term{normalization}{正规化}，$U'$ 为 $U$ 的一个\Term{open-subscheme}{开子概形}，在 $U_{\mathrm{red}}$ 的一维\Term{smooth-locus}{光滑轨迹}中\Term{dense}{稠密}，且 $G$ \Term{restriction}{限制}到 $U$ 后\Term{twisted-constant}{扭常值}：
\[
\begin{tikzcd}
U'\arrow[r,hook,"j'"]\arrow[d]&X'\arrow[d,"p"]&\\
U\arrow[r,hook,"j"]&X&Y\arrow[l,hook']
\end{tikzcd}
\]
则 $p_*j'_*G|Y\simeq j_*G|Y$。因而只须研究 $X$ \Term{smooth}{光滑}、一维且 $G|U$ \Term{twisted-constant}{扭常值}的情形。此时\Term{sheaf}{层} $G^\vee$ 为\Term{T-integral}{$T$-整的}，故已经知道，对每个 $x\in X$，$(j_*G^\vee)_{\bar x}$ 为\Term{T-integral}{$T$-整的}。
\SourcePage{399}
对 $x\in X$，用 $\mathfrak G_x$ 表示 $\widehat{\mathcal O}_{X,x}$ 的\Term{fraction-field}{分式域}的\Term{galois-group}{Galois 群}，用 $\mathfrak I_x$ 表示其\Term{inertia-group}{惯性群}。\Term{sheaf}{层} $G$ 给出 $\mathfrak G_x$ 的一个\Term{representation}{表示}，仍记为 $G$，而 $\mathfrak G_x/\mathfrak I_x=\operatorname{Gal}(\overline{k(x)}/k(x))$-模 $(j_*G^\vee)_{\bar x}$（相应地，$(j_*G)_{\bar x}^\vee$）便与 $G^\vee$ 中的\Term{invariant}{不变量}（相应地，\Term{coinvariant}{余不变量}）构成的 $\mathfrak G_x/\mathfrak I_x$-模等同。剩下只须证明下述\Term{lemma}{引理}。
\end{proof}

\begin{SGAresult}{\Term{lemma}{引理}}{5.3.1}
设 $K$ 为\Term{residue-field}{剩余域}为 $\mathbb F_q$ 的\Term{local-field}{局部域}，$G$ 为其\Term{galois-group}{Galois 群}，$I$ 为\Term{inertia-group}{惯性群}，$\rho$ 为 $G$ 在 $E$ 上 $n$ 维\Term{vector-space}{向量空间} $V$ 中的一个\Term{continuous-representation}{连续表示}（参见 5.0），$V^I$ 为 $I$ 下的\Term{invariant}{不变量}\Term{space}{空间}，$V_I$ 为 $I$ 下的\Term{coinvariant}{余不变量}\Term{space}{空间}，二者均看作 $G/I$-模。设 $\alpha_1,\ldots,\alpha_n$ 为 $F$ \Term{action}{作用}于 $V^I$ 的\Term{eigenvalue}{特征值}。则存在整数 $m_i\geq0$，使 $F$ \Term{action}{作用}于 $V_I$ 的\Term{eigenvalue}{特征值}为 $\alpha_iq^{m_i}$。
\end{SGAresult}
\begin{proof}
设 $I'$ 为 $I$ 中阶与 $\ell$ 互素的最大\Term{subgroup}{子群}。群 $\rho(I')$ 是\Term{finite-group}{有限群}，因为它是 $\operatorname{GL}(V)$ 的一个阶与 $\ell$ 互素的紧\Term{subgroup}{子群}。因而
\[
V_I\simeq\bigl(V^{\rho(I')}\bigr)_I.
\]
将 $V$ 换成 $V^{I'}$ 后，可以假设 $I'$ 在 $V$ 上的\Term{action}{作用}平凡。作为 $G/I$-模，有 $I/I'\simeq\mathbb Z_\ell(1)$；设 $T$ 为 $I/I'$ 的一个\Term{topological-generators}{拓扑生成元}，对每个 $\alpha$，以 $V^\alpha$ 表示 $T$ 在 $V$ 中对应于\Term{eigenvalue}{特征值} $\alpha$ 的\Term{generalized-eigenspace}{广义特征子空间}：有 $V=\bigoplus V^{(\alpha)}$；而\Term{resolution}{分解}
\[
V=V^{(1)}\oplus\bigoplus_{\alpha\ne1}V^{(\alpha)}
\]
在 $G$ 下\Term{stable}{稳定}，且与取\Term{dual}{对偶}相容。
\SourcePage{400}
因而可以将 $V$ 换成 $V^{(1)}$，从而假定 $I$ 的\Term{action}{作用}是\Term{unipotent}{幺幂}的。
\end{proof}

\begin{SGAresult}{\Term{lemma}{引理}}{5.3.2}
设 $V$ 为\Term{characteristic-zero}{特征零}的域 $k$ 上的\Term{finite-dimensional-vector-space}{有限维向量空间}，$I$ 为一个单参数\Term{additive-group}{加法群}，$\rho:I\to\operatorname{GL}(V)$ 为 $I$ 在 $V$ 中的一个\Term{representation}{表示}，$N$ 为 $I$ 的一个\Term{infinitesimal-generator}{无穷小生成元}，并置
\[
V^{I,k}=V^I\cap\operatorname{Im}(N^k),\qquad
V_{I,k}=\operatorname{Im}\bigl(\operatorname{Ker}(N^{k+1})\longrightarrow V_I\bigr).
\]
则映射 $N^k$ 诱导一个\Term{isomorphism}{同构}
\[
N^k:V_{I,k}/V_{I,k-1}\xrightarrow{\sim}V^{I,k}/V^{I,k+1},
\]
并定义\Term{canonical}{典范}\Term{isomorphism}{同构}
\[
V_{I,k}/V_{I,k-1}\otimes I^{\otimes k}\simeq V^{I,k}/V^{I,k+1}.
\]
\end{SGAresult}
\begin{proof}
通过 $N$，$I$ 的\Term{representation}{表示} $V$ 与被 $T$ 的某个\Term{nilpotent}{幂零}化的 $k[T]$-模等同。只须在 $V$ 对应于一个\Term{irreducible}{不可约} $k[T]$-模 $k[T]/T^k$ 时验证 3.2，而这是平凡的。

现完成 5.3.1 的证明。5.3.2 这一组结论在 $V_I$ 和 $V^I$ 上分别给出递增与递减\Term{filtration}{滤过} $V_{I,k}$ 与 $V^{I,k}$，以及 $G/I$-模的\Term{isomorphism}{同构}
\[
V^{I,k}/V^{I,k+1}\simeq V_{I,k}/V_{I,k-1}(k),
\]
断言随之得证。
\end{proof}

\SourcePage{401}
\begin{SGAresult}{\Term{theorem}{定理}}{5.4}
设 $X$ 为 $\mathbb F_q$ 上\Term{dimension}{维数}为 $n$ 的\Term{separated}{分离}\Term{finite-type}{有限型}\Term{scheme}{概形}，$G$ 为 $X$ 上的\Term{constructible}{可构造} $E$-\Term{sheaf}{层}。假定对每个 $x\in X$，$G_{\bar x}^\vee$ 的\Term{endomorphism}{自同态} $F_x$ 均为\Term{T-integral}{$T$-整的}。则 $\underline H_c^i(X,G)^\vee$ 的\Term{endomorphism}{自同态} $q^iF^\vee$ 为\Term{T-integral}{$T$-整的}，其\Term{endomorphism}{自同态} $q^nF^\vee$ 也为\Term{T-integral}{$T$-整的}。
\end{SGAresult}
\begin{proof}
像\Term{theorem}{定理} 5.2 的证明那样，归约到 $X$ 为\Term{smooth}{光滑}\Term{curve}{曲线}、$G$ 在 $X$ 上\Term{twisted-constant}{扭常值}的情形。设 $j$ 为将 $X$ \Term{embedding}{嵌入}其\Term{compactification}{紧化}、\Term{projective}{射影}\Term{smooth}{光滑}\Term{curve}{曲线} $X'$ 的\Term{inclusion}{包含映射}。

由\Term{poincare-duality}{Poincaré 对偶}，$\underline H^i(X',j_*G)$ 与 $\underline H^{2-i}(X',j_*G^\vee(1))$ 互为\Term{dual}{对偶}。由推论 5.3 (a)，\Term{theorem}{定理} 5.2 适用于 $j_*G^\vee$：
\[
\begin{aligned}
\underline H^0(X',j_*G^\vee)&\simeq\underline H^2(X',j_*G)^\vee(-1)
&&\text{为 $T$-整的},\\
\underline H^1(X',j_*G^\vee)&\simeq\underline H^1(X',j_*G)^\vee(-1)
&&\text{为 $T$-整的},\\
\underline H^2(X',j_*G^\vee)(1)&\simeq\underline H^0(X',j_*G)^\vee
&&\text{为 $T$-整的}.
\end{aligned}
\]
这就对 $j_*G$ 证明了\Term{theorem}{定理}。以\Term{exact-sequence}{正合列}
\[
0\longrightarrow j_!G\longrightarrow j_*G\longrightarrow H\longrightarrow0
\]
定义 $H$。若 $Y$ 为 $H$ 的\Term{support}{支撑}，则\Term{cohomology}{上同调}\Term{exact-sequence}{正合列}给出
\[
\begin{gathered}
\underline H_c^0(X,G)\hookrightarrow\underline H^0(X',j_*G),\\
\underline H^0(Y,H)\longrightarrow\underline H_c^1(X,G)\longrightarrow\underline H^1(X',j_*G),\\
\underline H_c^2(X,G)\xrightarrow{\sim}\underline H^2(X',j_*G),
\end{gathered}
\]
\SourcePage{402}
故 $(X,G)$ 的\Term{theorem}{定理}由 $(X',j_*G)$ 的\Term{theorem}{定理}\Term{pushout}{推出}。
\end{proof}

\textbf{5.5。}上述\Term{fibration}{纤维化}为\Term{curve}{曲线}的方法，也可给出 Frobenius \Term{eigenvalue}{特征值}的复绝对值的 Lang--Weil 型上界 [3]。

设 $X$ 为 $\mathbb F_q$ 上的\Term{finite-type}{有限型}\Term{scheme}{概形}，$G$ 为 $X$ 上的\Term{constructible}{可构造} $E$-\Term{sheaf}{层}，$\lambda\in\mathbb R$。若对每个 $x\in X$，置 $\#k(x)=q_x$，则 $F_x$ \Term{action}{作用}于 $G_{\bar x}$ 的\Term{eigenvalue}{特征值}都是\Term{algebra}{代数}数 $\alpha$，其所有复绝对值均满足
\[
|\alpha|\leq q_x^\lambda,
\]
便记作
\[
\log|G|\leq\lambda.
\]

\begin{SGAresult}{\Term{lemma}{引理}}{5.5.1}
设 $X$ 为 $\mathbb F_q$ 上的一条\Term{curve}{曲线}，这里是指 $\mathbb F_q$ 上\Term{dimension}{维数} $\leq1$ 的\Term{separated}{分离}\Term{finite-type}{有限型}\Term{scheme}{概形}。设 $G$ 为 $X$ 上的\Term{constructible}{可构造} $E$-\Term{sheaf}{层}，且 $\log|G|\leq\lambda$，其中 $\lambda\in\mathbb R$。则 $F$ \Term{action}{作用}于 $\underline H_c^1(X,G)$ 的\Term{eigenvalue}{特征值} $\alpha$ 的复绝对值满足 $|\alpha|\leq q^{\lambda+1}$。
\end{SGAresult}
\begin{proof}
设 $\sigma$ 为 $E$ 的一个复\Term{embedding}{嵌入}，或只取 $E$ 中 $\mathbb Q$ 的\Term{algebraic-closure}{代数闭包}的一个复\Term{embedding}{嵌入}。

只须验证：无论 $\sigma$ 如何选取，\Term{rational-function}{有理函数} $\sigma Z(X,G,t)$ 的零点 $\beta$ 均满足 $|\beta^{-1}|\leq q^{\lambda+1}$。
\SourcePage{403}
事实上，由 5.2.1 已知，若 $\alpha$ 是 Frobenius \Term{action}{作用}于 $\underline H_c^0(X,G)$ 或 $\underline H_c^2(X,G)$ 的一个\Term{eigenvalue}{特征值}，则 $|\sigma\alpha|\leq q^{\lambda+1}$。

来验证级数
\[
\log\sigma Z(X,G,t)
=\sum_k\frac{t^k}{k}\sum_{x\in X(\mathbb F_{q^k})}\sigma\operatorname{Tr}(F,x^*G)
\]
当 $|t|<q^{-\lambda-1}$ 时收敛。将 $X$ 表示成 $\mathbb P^1$ 的\Term{ramified-covering}{分歧覆盖}，便立即得到上界
\[
\#X(\mathbb F_{q^k})\leq A q^k.
\]
另一方面，由假设有
\[
|\sigma\operatorname{Tr}(F,x^*G)|\leq B(q^k)^\lambda.
\]
所以 $\sigma Z(X,G,t)$ 中 $t^k$ 的系数 $c_k$ 满足
\[
|c_k|\leq\frac1k\,C\,(q^{1+\lambda})^k,
\]
断言随之成立。
\end{proof}

\begin{SGAresult}{\Term{theorem}{定理}}{5.5.2}
设 $X$ 为 $\mathbb F_q$ 上\Term{dimension}{维数} $\leq n$ 的\Term{separated}{分离}\Term{finite-type}{有限型}\Term{scheme}{概形}，$G$ 为 $X$ 上的\Term{constructible}{可构造} $E$-\Term{sheaf}{层}，且 $\log|G|\leq\lambda$，其中 $\lambda\in\mathbb R$。

则 Frobenius \Term{action}{作用}于 $\underline H_c^i(X,G)$ 的\Term{eigenvalue}{特征值} $\alpha$ 的复绝对值满足
\[
|\alpha|\leq q^{\lambda+i}\qquad\text{且}\qquad |\alpha|\leq q^{\lambda+n}.
\]
\end{SGAresult}
\SourcePage{404}
\begin{proof}
若 $n\leq0$，5.5.2 是平凡的。若 $n=1$，5.5.2 中 $i=0$ 与 $i=2$ 的情形由 5.2.1 得到，而 $i=1$ 的情形由 5.5.1 得到。然后对 $n$ 归纳，如同 5.5.2 的证明。
\end{proof}

\begin{SGAresult}{推论}{5.5.3}
设 $X$ 为 $\mathbb F_q$ 上\Term{dimension}{维数}为 $n$ 的\Term{separated}{分离}\Term{finite-type}{有限型}\Term{scheme}{概形}。则几何 Frobenius \Term{endomorphism}{自同态} $F$ \Term{action}{作用}于 $\underline H_c^i(X,\mathbb Q_\ell)$ 的\Term{eigenvalue}{特征值} $\alpha$ 都是\Term{algebraic-integer}{代数整数}。此外：

\textup{(i)} 对 $\ell'\ne p$，$\alpha$ 是一个 $\ell'$-进\Term{unit}{单位}。

\textup{(ii)} 当 $i>n$ 时，$\alpha$ 被 $q^{i-n}$ 整除。

\textup{(iii)} 当 $i\leq n$ 时（相应地，当 $i\geq n$ 时），$q^i\alpha^{-1}$（相应地，$q^n\alpha^{-1}$）是\Term{algebraic-integer}{代数整数}。

\textup{(iv)} 当 $0<i$ 时，$\alpha$ 的复绝对值满足
\[
|\alpha|\leq q^{i-\frac12}\qquad\text{且}\qquad |\alpha|\leq q^n.
\]
\end{SGAresult}
\begin{proof}
由 5.2.2，$\alpha$ 与 $q^{n-i}\alpha$ 均为\Term{algebraic-integer}{代数整数}。断言 (i) 至 (iii) 随即由 5.4 \Term{pushout}{推出}。对 $i>n$，(iv) 由 5.5.2 得到。为在 $i\leq n$ 时得到 (iv)，须将 5.5.2 与\Term{curve}{曲线}的 Riemann 假设结合。对 $n$ 归纳，$n=0$ 的情形是平凡的。若 $U$ 是 $X$ 的一个\Term{open-subscheme}{开子概形}，其补集 $Z$ 的\Term{dimension}{维数} $<n$，则\Term{exact-sequence}{正合列}
\[
\underline H_c^i(U,\mathbb Q_\ell)\longrightarrow
\underline H_c^i(X,\mathbb Q_\ell)\longrightarrow
\underline H_c^i(Z,\mathbb Q_\ell)
\]
表明只须对 $U$ 证明 5.2.2。预先将 $X$ 换成 $X_{\mathrm{red}}$ 后，便归约到 $X$ 为\Term{smooth}{光滑}\Term{affine}{仿射}的情形，
\SourcePage{405}
而且存在一个\Term{smooth-morphism}{光滑态射} $f:X\to S$ 满足下述条件：

\textup{(a)} $\dim S\leq n-1$；

\textup{(b)} $f$ 的\Term{fibre}{纤维}均为纯一维的。

由于 $f$ 的\Term{fibre}{纤维}都是\Term{affine}{仿射}\Term{curve}{曲线}，有
\[
R^0f_!\mathbb Q_\ell=0.
\]
由 Weil [4]，有
\[
\log|R^1f_!\mathbb Q_\ell|\leq\tfrac12,
\tag{5.5.3.1}
\]
将 5.2.1 用于 $f$ 的\Term{fibre}{纤维}，又有
\[
\log|R^2f_!\mathbb Q_\ell|\leq1.
\tag{5.5.3.2}
\]
这里，$f$ 的 Leray \Term{spectral-sequence}{谱序列}化为一个\Term{long-exact-sequence}{长正合列}
\[
\underline H_c^{i-2}(S,R^2f_!\mathbb Q_\ell)\longrightarrow
\underline H_c^i(X,\mathbb Q_\ell)\longrightarrow
\underline H_c^{i-1}(S,R^1f_!\mathbb Q_\ell).
\]
由 (5.5.3.1)（相应地，(5.5.3.2)）及 5.5.2，Frobenius \Term{action}{作用}于 $\underline H_c^{i-1}(S,R^1f_!\mathbb Q_\ell)$（相应地，$\underline H_c^{i-2}(S,R^2f_!\mathbb Q_\ell)$）的\Term{eigenvalue}{特征值} $\alpha$ 的复绝对值满足
\[
|\alpha|\leq q^{i-1+\frac12}=q^{i-\frac12}
\]
（相应地，$|\alpha|\leq q^{i-2+1}=q^{i-1}\leq q^{i-\frac12}$）。推论随之成立。
\end{proof}

\SourcePage{406}
为了在无\Term{support}{支撑}的\Term{cohomology}{上同调}中得到一些类似上述结果的结论，似乎必须有\Term{resolution-of-singularities}{奇点消解}。\Term{dimension}{维数} $\leq1$ 时，推论 5.3 已给出这一问题的解答。

\begin{SGAresult}{\Term{theorem}{定理}}{5.6}
设 $N$ 为整数，并假定对于 $\mathbb F_q$ 上\Term{dimension}{维数} $\leq N$ 的\Term{finite-type}{有限型}\Term{scheme}{概形}，存在 Hironaka 强\Term{form}{形式}的\Term{resolution-of-singularities}{奇点消解}（截至目前，可取 $N=2$ [1]）。设 $f:X\to Y$ 为 $\mathbb F_q$ 上\Term{finite-type}{有限型}\Term{scheme}{概形}之间的\Term{morphism}{态射}，$G$ 为 $X$ 上的\Term{T-integral}{$T$-整}\Term{constructible}{可构造} $E$-\Term{sheaf}{层}。则若 $\dim X\leq N$，$Y$ 上的 $E$-\Term{sheaf}{层} $R^if_*G$ 均为\Term{T-integral}{$T$-整的}。
\end{SGAresult}
\begin{proof}
证明对 $X$ 的\Term{dimension}{维数} $n\leq N$ 归纳。设 $j:U\to X$ 为 $X$ 的一个\Term{smooth}{光滑}\Term{affine}{仿射}\Term{open-subscheme}{开子概形}，使 $G$ \Term{restriction}{限制}到 $U$ 后\Term{twisted-constant}{扭常值}，并使其补集 $Z$ 的\Term{dimension}{维数} $<n$。利用\Term{resolution-of-singularities}{奇点消解}假设，知道 $R^ij_*(j^*G)$ 均\Term{constructible}{可构造}（SGA 5 第 I 讲\SGAsourcegap{原书此处在 SGA 5 I 后留白，未填具体引注号。}，[2]）；当 $i>0$ 时，它们集中在 $Z$ 上，而当 $i=0$ 时，$j_*(j^*G)$ 只在 $Z$ 上与 $G$ 不同。因此，由归纳假设，只须对\Term{sheaf}{层} $j^*G$ 及\Term{morphism}{态射} $j$、$fj$ 验证\Term{theorem}{定理}。故只须在额外假设 $X$ \Term{affine}{仿射}\Term{smooth}{光滑}、$G$ \Term{twisted-constant}{扭常值}的条件下验证\Term{theorem}{定理}（即 $G$ 由 $X$ 的\Term{fundamental-group}{基本群}在一个\Term{finite-dimensional}{有限维} $E$-\Term{vector-space}{向量空间}中的\Term{continuous-representation}{连续表示}给出）。

\SourcePage{407}
设 $R$ 为 $E$ 的\Term{ring-of-integers}{整数环}，$\lambda$ 为一个\Term{uniformizer}{一致化参数}，$e$ 为绝对\Term{ramification-index}{分歧指数}，$b$ 为一个 $>e/(\ell-1)$ 的整数。由于 $X$ \Term{smooth}{光滑}，故为\Term{normal}{正规}且\Term{fundamental-group}{基本群}为\Term{profinite}{投有限}的，存在一个\Term{twisted-constant}{扭常值} $R$-\Term{sheaf}{层} $G^0$，使 $G\simeq G^0\otimes_RE$。设 $p:X'\to X$ 为使\Term{locally-constant}{局部常值}\Term{sheaf}{层} $G^0\otimes_RR/\lambda^b$ \Term{trivialization}{平凡化}的\Term{etale-covering}{平展覆盖}。\Term{sheaf}{层} $G$ 是 $p_*p^*G$ 的一个\Term{direct-factor}{直因子}，故只须对 $p^*G$ 与 $f\circ p$ 验证\Term{theorem}{定理}。因此，只须处理满足下述假设的情形：

\textup{(*)} $X$ \Term{affine}{仿射}且\Term{smooth}{光滑}，并存在一个\Term{twisted-constant}{扭常值} $R$-\Term{sheaf}{层} $G^0$，使 $G\simeq G^0\otimes_RE$，且\Term{sheaf}{层} $G^0\otimes_RR/\lambda^b$ 为常值的。

将\Term{morphism}{态射} $f$ \Term{resolution}{分解}为一个\Term{open-immersion}{开浸入}与一个\Term{proper}{固有}\Term{morphism}{态射}的复合 $f=\bar f j$：
\[
\begin{tikzcd}
X\arrow[r,hook,"j"]\arrow[d,"f"']&X'\arrow[dl,"\bar f"]\\
Y&
\end{tikzcd}
\]
\Term{theorem}{定理}已对 $\bar f$ 证明（5.2），故只须对 $j$ 证明。将\Term{resolution-of-singularities}{奇点消解}用于 $X'$，可假设 $X'$ \Term{smooth}{光滑}，且 $X$ 为 $X'$ 中一个\Term{normal-crossings-divisor}{正规交叉除子}的补集，该\Term{divisor}{除子}是若干\Term{smooth}{光滑}\Term{divisor}{除子}的并。假设 (*) 蕴含\Term{sheaf}{层} $G$ 沿此\Term{divisor}{除子}\Term{tamely-ramified}{温和分歧}。此时有下述\Term{lemma}{引理}。
\end{proof}

\SourcePage{408}
\begin{SGAresult}{\Term{lemma}{引理}}{5.6.1}
设 $X$ 为 $\mathbb F_q$ 上的\Term{smooth}{光滑}\Term{finite-type}{有限型}\Term{scheme}{概形}，$D$ 为 $X$ 中一个\Term{normal-crossings-divisor}{正规交叉除子}，且是若干\Term{smooth}{光滑}\Term{divisor}{除子} $D_i$ 的并，$j:U\hookrightarrow X$ 为 $D$ 在 $X$ 中的补集，$G$ 为 $U$ 上一个沿 $D$ \Term{tamely-ramified}{温和分歧}的\Term{twisted-constant}{扭常值} $E$-\Term{sheaf}{层}。则若 $G$ 为\Term{T-integral}{$T$-整的}，各\Term{sheaf}{层} $R^ij_*G$ 均为\Term{T-integral}{$T$-整的}。
\end{SGAresult}

我们承认下述\Term{lemma}{引理}（参见 SGA 1 第 XIII 讲\SGAsourcegap{原书此处在 SGA 1 XIII 后留白，未填具体引注号。}）：

\begin{SGAresult}{\Term{lemma}{引理}}{5.6.2}
\textup{(i)} 在 5.6.1 的假设下，设 $D'$ 为 $D$ 的一个\Term{component}{分支}，$D''$ 为其余\Term{component}{分支}的并，$U'=X-D''$，$D'^0=D'\cap U'$：
\[
\begin{tikzcd}
&D'^0\arrow[r,hook,"j'_D"]\arrow[d,"i^0"']&D'\arrow[d,"i"]\\
U\arrow[r,hook,"j_1"]&U'\arrow[r,hook,"j'"]&X.
\end{tikzcd}
\]
则\Term{base-change-morphism}{基变换映射}
\[
i^*Rj'_*(Rj_{1*}G)\longrightarrow Rj'_{D*}\,i^{0*}(Rj_{1*}G)
\]
为\Term{isomorphism}{同构}，而且各\Term{sheaf}{层} $i^{0*}R^ij_{1*}G$ 均\Term{twisted-constant}{扭常值}，并沿 $D'-D'^0$ \Term{tamely-ramified}{温和分歧}。

\textup{(ii)} 若 $C$ 是 $X$ 上的一条\Term{smooth}{光滑}\Term{curve}{曲线}，在 $D$ 的一个\Term{smooth-point}{光滑点} $x$ 处与 $D$ \Term{transverse}{横截}，
\[
\begin{tikzcd}
C\cap U\arrow[r,hook,"j_C"]\arrow[d,"k'"']&C\arrow[d,"k"]\\
U\arrow[r,hook,"j"]&X,
\end{tikzcd}
\]
\SourcePage{409}
则\Term{base-change-morphism}{基变换映射}
\[
k^*Rj_*G\longrightarrow Rj_{C*}\,k'^*G
\]
在 $x$ 处是\Term{isomorphism}{同构}。
\end{SGAresult}
\begin{proof}[\Term{lemma}{引理} 5.6.1 的证明]
对\Term{dimension}{维数}归纳，并暂且承认 5.6.1 对\Term{curve}{曲线}成立，由此从 5.6.2 \Term{pushout}{推出} 5.6.1。由 5.6.2 (ii) 和这一假定，\Term{sheaf}{层} $R^ij_*G$ 在 $D$ 的每个\Term{smooth-point}{光滑点}处均为\Term{T-integral}{$T$-整的}。特别地，采用 5.6.2 (i) 的记号，各\Term{sheaf}{层} $R^ij_{1*}G$ 为\Term{T-integral}{$T$-整的}，因而由 6.2 (i)，并将归纳假设用于 $i^{0*}R^ij_{1*}G$ 与 $j'_D$，得知无论选取 $D$ 的哪一个\Term{component}{分支} $D'$，各\Term{sheaf}{层} $i^*R^ij_*G$ 均为\Term{T-integral}{$T$-整的}。这就证明了 5.6.1 的结论。

剩下的是\Term{curve}{曲线}的情形。对一条满足 5.6.1 假设的\Term{curve}{曲线}，若 $x$ 为 $X-U$ 的一点，$\mathfrak G$ 为 $\widehat{\mathcal O}_{x,X}$ 的\Term{fraction-field}{分式域}的\Term{galois-group}{Galois 群}，$\mathfrak I$ 为\Term{inertia-group}{惯性群}，仍用 $G$ 表示\Term{sheaf}{层} $G$ 定义的 $\mathfrak G$ 的\Term{representation}{表示}，并设 $\bar x$ 为 $X$ 的一个位于 $x$ 之上的\Term{geometric-point}{几何点}，则有\Term{galois-module}{Galois 模}（即 $\mathfrak G/\mathfrak I$-模）的\Term{isomorphism}{同构}：
\[
\begin{aligned}
(j_*G)_{\bar x}&\simeq G^{\mathfrak I},\\
(R^1j_*G)_{\bar x}&\simeq G_{\mathfrak I}(-1),\\
(R^ij_*G)_{\bar x}&=0\qquad\text{当 }i\geq2.
\end{aligned}
\]
断言由 5.3 (a) 及 5.3.1 得到。
\end{proof}

\SourcePage{410}
\textbf{参考文献}

\begin{enumerate}
\item[{[1]}] S. S. ABHYANKAR，\Term{embedding}{嵌入}\Term{algebra}{代数}\Term{surface}{曲面}的\Term{resolution-singularities}{奇点消解}（\textit{Resolution of singularities of embedded algebraic surfaces}），Acad. Press，New York，1966。
\item[{[2]}] J. P. JOUANOLOU，\Term{scheme}{概形}的 $\ell$-进\Term{cohomology}{上同调}（\textit{Cohomologie $\ell$-adique des schémas}），将刊于 \textit{Lecture Notes of Math.}。
\item[{[3]}] S. LANG 与 A. WEIL，\Term{finite-field}{有限域}上簇的点数（\textit{Number of points of varieties in finite fields}），\textit{Am. J. Math.} \textbf{76}（1954），819--827。
\item[{[4]}] A. WEIL，论\Term{algebraic-curve}{代数曲线}及由其\Term{derive}{导出}的簇（\textit{Sur les courbes algébriques et les variétés qui s'en déduisent}），\textit{Act. Sci. Ind.} 1041，Hermann，1948。
\end{enumerate}
